\section{Structural Timing Bias and ECC Removal}
\label{sec:trng-structural}

\subsection{Experimental Platform}
The experiments in this chapter use the same experimental platform
described in Section~\ref{sec:brewrc-setup}: two MB85AS4MT (4M) and two
MB85AS8MT (8M) standalone ReRAM modules connected to a DE10-Nano FPGA
development kit via SPI. Write latency measurements are captured using
the same hardware timer and WIP polling methodology. The ECC parity
submatrices $P_{4M}$ and $P_{8M}$ recovered in Chapter~\ref{chap:brewrc}
are used throughout.

\subsection{Sierpinski Structure and ECC Activity}

\begin{figure}[t!]
	\centering
	\includegraphics[width=\linewidth]{img/single-byte-sierpinski-heatmaps.pdf}
	\caption{
		Single-byte transition Sierpinski experiment showing a heatmap of single-byte transition latencies.
		Three 4M and three 8M chips are initialized to a global bit state, G=$\mathbf{0}$ or G=$\mathbf{1}$. Two randomly selected bytes (Byte 1, and 0) is exhaustively transitioned, FROM an initial value, TO every possible value. Each pixel is colored according to the median write latency.
		The underlying Sierpinski-like pattern replicates the result on 4M in \cite{tolga2024trng} and extends it to 8M.
		\label{fig:sierp-orig}
	}
\end{figure}

\begin{figure}[t!]
	\centering
	\includegraphics[width=\linewidth]{img/single-byte-sierpinski-heatmaps-resid.pdf}
	\caption{
		Single-byte, transition experiment showing a heatmap of byte transitions latencies, with an ECC-aware latency model subtracted.
		We construct an ECC-aware timing model using techniques developed in Section~\ref{sec:brewrc-method}, and subtract it from the median latencies collected in Figure~\ref{fig:sierp-orig}, $\hat{t}_{ws}-t_w$.
		The structural timing bias introduced by ECC are fully removed, leaving a noisy latency component.
		This final noise can be profiled, and optimized to harvest entropy.
		\label{fig:sierp-resid}
	}
\end{figure}

Arul et al.~\cite{tolga2024trng} identified a Sierpinski-like pattern in
single-byte transition latency heatmaps of the MB85AS4MT and attributed
it to the bit-level polarity structure of data transitions. The
experiments in Chapter~\ref{chap:brewrc} (Section~\ref{sec:brewrc-exp4})
reproduce this pattern and demonstrate that it arises from ECC parity
switching activity rather than data polarity alone. This section
confirms this attribution by constructing an ECC-aware timing model and
subtracting its predictions from measured latencies. Each chip is
initialized to a global bit state, $G=\mathbf{0}$ or $G=\mathbf{1}$.
Each byte under test is exhaustively transitioned from every initial
value to every final value. Each FROM/TO transition is repeated 10 times
and the experiment is performed on all byte positions in 10 randomly
selected words, on three modules each of 4M and 8M density.

Figure~\ref{fig:sierpinski} (top) reproduces the Sierpinski-like pattern
on both 4M and 8M devices, extending the result of~\cite{tolga2024trng}
to the 8M for the first time. The ECC-aware timing model
$\hat{t}_w(\delta)$ (Equation~\ref{eq:timing-model}) is then used to
predict the median latency for each byte transition by computing the
codeword DHD using the recovered ECC. Figure~\ref{fig:sierpinski}
(bottom) shows the residual heatmap after subtracting the predicted
latency from the measured median latency. The structural pattern is
completely eliminated, leaving only a noisy residual consistent with
stochastic filament dynamics. This result confirms that ECC activity is
the sole source of structured timing variance in these devices, and that
TRNGs that do not account for the underlying ECC will overstate their
harvested entropy.

